\subsubsection{Machine Learning}

Segundo \citeonline{SilvaVilela2021}, Machine Learning pode ser 
definida como um ramo da \gls{ia} responsável por capacitar computadores
à simular aprendizado com base em processamento de dados. 
Esta técnica está em ascensão atualmente, sendo amplamente aplicado em softwares que 
tem por objetivo a identificação de padrões, como é o caso do \gls{ocr},
detecção de phishing em e-mails, detecção de uso fraudulento de cartões de crédito, 
entre outras aplicações.

Tal subconjunto, tem por característica a distinção de métodos de aprendizagem, 
tendo como os mais difundidos os métodos supervisionado, 
não supervisionado e semi-supervisionado.

\subsubsubsection{Aprendizado Supervisionado}

Neste método de aprendizagem, o objetivo é ensinar ao modelo de \gls{ia}, 
através de objetos previamente rotulados, à identificar características de modo à 
classificar novos objetos, \cite{SilvaVilela2021}. 
Entre os algoritmos mais importantes para a aplicação deste tipo de 
método de treinamento estão:

\begin{enumerate}
	\item k-Nearest Neighbours;
    \item Regressão Linear;
    \item Máquina de Vetores de Suporte (SVM);
    \item Regressão Linear;
    \item Árvore de Decisão e Floresta Aleatória;
    \item Redes Neurais;
\end{enumerate}

Observação, estes algoritmos não necessariamente são concorrentes, 
portanto em único modelo pode haver implementação de um ou mais 
dos algoritmos citados, \cite{SilvaVilela2021}.

\subsubsection{Deep Learning}

Deep Learning, ou aprendizado profundo, é uma técnica que se encontra 
como em um subcampo do Machine Learning, onde seu desenvolvimento foi pautado 
na simulação do processo de aprendizagem humano. 
No caso utilizando como base o algoritmo de redes neurais, 
construindo estruturas conectadas, de forma à emular o comportamento do cérebro humano, 
\cite{jesus2023inteligencia}.


\subsubsection{Redes Neurais}

Como apresentado por \citeonline{jesus2023inteligencia}, 
redes neurais são estruturas baseadas em grafos, que são representadas em camadas, 
cada camada é compostas por nós, denominado neurônios, e arestas, 
denominada sinapses. a estrutura básica de uma rede neural possui 3 camadas: 
entrada, camada intermediária, saída, onde, respectivamente, são responsáveis 
por executar as funções de receber os dados externos, 
processar os dados e entregar o resultado.

A técnica de redes neurais apresenta derivações amplamente utilizadas como alicerce 
no desenvolvimento de sistemas baseados em Deep Learning.

\subsubsubsection{Redes Neurais Convolucionais}

As \gls{cnn}, são amplamente aplicadas em processamento de dados não estruturados, 
como imagens, vídeos e áudios. Tal técnica tem como característica o 
processamento de pequenos trechos de dados por vez, ao contrário 
de outras técnicas que tentam realizar o processamento de forma atômica, 
\cite{jesus2023inteligencia}.

Como apresentado por \citeonline{Atika2022ImageSegmentation}, 
o desenvolvimento de uma CNN será apoiada na implementação de uma arquitetura, 
tal técnica abrange uma vasta quantidade de arquiteturas que foram 
desenvolvidas para auxiliar e alcançar melhores resultados à 
determinadas aplicações do software.
Dado o objetivo da implementação de CNN neste trabalho, 
destaca-se a arquitetura U-Net. 

Como \citeonline{Atika2022ImageSegmentation} explica,
a arquitetura U-Net tem por objetivo a definição de bordas dentro de imagens, 
onde é amplamente utilizado na biomedicina com o foco de delimitar 
estruturas presentes em imagens, tal característica justifica a aplicação 
deste método com o objetivo de definir \gls{roi}. 
A arquitetura U-Net é caracterizada pela presença de três etapas de 
processamento: Contração, Gargalo e Expansão. 
A camada de contração é composta por vários blocos, 
onde cada bloco recebe uma entrada e aplica uma camada de convolução seguida por 
uma camada de operação de max-pooling, cada uma das camadas são matrizes, 
que respectivamente, apresentam a estrutura 3x3 e 2x2.

Um adendo, as camadas que constroem a arquitetura U-Net utilizam função de ativação ReLu 
apresentada na equação~\ref{eq:relu_def}.

\begin{equation}
    f(x) = \text{ReLU}(x) = \max(0, x) = 
    \begin{cases} 
        x, & \text{se } x > 0 \\ 
        0, & \text{se } x \le 0 
    \end{cases}
    \label{eq:relu_def}
\end{equation}

\subsubsection{Quantização e QAT}

Quantização é um método de simplificação dos parâmetros de um modelo de \gls{rna}. 
Como o \citeonline{ferreira2022compressao} explica, a aplicação de tal técnica, 
traz como benefício a redução do tamanho de modelos, 
viabilizando a utilização em equipamentos de baixo poder computacional, 
além do aumento de velocidade de inferência e simplificação de cálculos realizados pelo processador, 
frente à um modelo que na maioria dos casos possuem seus parâmetros representados em pontos 
flutuantes de 32 bits. 

\gls{qat} é a aplicação da técnica de quantização pré-treinamento, 
treinando o modelo à reconhecer padrões em imagens utilizando parâmetros de baixa precisão.

\subsubsection{OCR - Fundamentos e Escopo}
O \gls{ocr} é um subcampo da Visão Computacional e da Inteligência Artificial 
dedicado à conversão de imagens contendo texto (impresso, manuscrito ou gravado) 
em dados legíveis por máquina \cite{chen2021survey}. 
Originalmente voltado à digitalização de documentos, 
o escopo foi ampliado pelo \textit{Deep Learning} 
para aplicações industriais e de IoT, \cite{jawas2018water}, 

\subsubsubsection{Pipeline Tradicional}
Um sistema clássico de OCR opera em etapas interdependentes \cite{chen2021survey}:
\begin{enumerate}
    \item \textbf{Aquisição e pré-processamento}: melhoria de contraste, correção de inclinação (\textit{deskewing}), remoção de ruído e binarização; em dispositivos embarcados, etapas de segmentação semântica podem isolar a Região de Interesse (ROI) antes do reconhecimento.
    \item \textbf{Localização e segmentação de texto}: identificação das regiões de texto.
    \item \textbf{Reconhecimento de caracteres}: classificação dos segmentos ou da sequência completa.
    \item \textbf{Pós-processamento}: dicionários, correção ortográfica ou verificação sintática para mitigar erros residuais.
\end{enumerate}

\subsubsubsection{Evolução Tecnológica: dos Métodos Heurísticos ao Deep Learning}
A literatura descreve três gerações principais \cite{chen2021survey}:

\subparagraph{Primeira geração (Template Matching).}
Sobreposição pixel a pixel com gabaritos estáticos; 
alta vulnerabilidade a variações de escala, rotação e fontes.

\subparagraph{Segunda geração (Características manuais + classificadores clássicos).}
Uso de descritores como HOG e SIFT, 
combinados com SVM ou MLP; dependência de \textit{feature engineering}.

\subparagraph{Terceira geração (End-to-End Deep Learning).}
Elimina a segmentação explícita de caracteres; 
modelos \textit{end-to-end} aprendem simultaneamente representações 
visuais e dependência contextual sequencial \cite{shi2017crnn}.

\subsubsection{Arquitetura CRNN e CTC}
A \gls{ctc}, é uma técnica utilizada para 
inferência de classificação em sequências de 
comprimento variável (palavras, placas, valores, entre outros),
que introduz um \textit{token} oculto de espaçamento (\textit{blank}) 
e marginaliza todos os alinhamentos temporais válidos, 
permitindo dotar modelos com a capacidade de inferir sequências de classes, ao mesmo tempo que
evita o colapso de repetições consecutivas (ex.: dígitos idênticos adjacentes) \cite{shi2017crnn, graves2006ctc}.
 
a arquitetura \gls{crnn}, é amplamente utilizada em aplicações que necessitam de realizar
inferência de classes em sequência, tal arquitetura combina a arquitetura \gls{cnn} e 
\gls{rnn} com a perda \gls{ctc}.